\SourceLeaf{169}die Summe aller Annäherungen, so oft man sie auch fortsetzte,
doch immer kleiner als $\frac{1}{1-\psi} \times ckf$ bleiben, und
diese Grenze könnte dann immer noch kleiner als der rechte Winkel
$dk\varphi$ sein. Du hast mein aufrichtiges Urteil verlangt. Ich
habe es gegeben und versichere Dir noch einmal, dass
es mich innig freuen wird, wenn Du alle Schwierigkeiten überwindest.

.\, .\, .\, .\, .\, .\, .\, .\, .\, .\, .\, .\, .\, .\, .\, .\, .

\medskip
\begin{center}
\rule{0.3\textwidth}{0.4pt}
\end{center}
\medskip

\begin{center}
\textsc{Bemerkungen.}
\end{center}

\textsc{Gauss} und \textsc{Wolfgang Bolyai} aus Bolya in
Siebenbürgen haben vom Herbst 1796 bis Herbst 1798 zusammen in
Göttingen studiert und freundschaftlich miteinander verkehrt. Vor
\textsc{Bolyais} Heimreise sind sie noch einmal am 25. Mai 1799 in
Clausthal zusammengetroffen. Der Brief vom 16. Dezember 1799, von
dem hier ein Auszug mitgeteilt wird, ist der erste, den \textsc{Gauss}
nach dem Abschied an seinen Freund gerichtet hat.

\textsc{Bolyai} hat seine »Göttingische Theorie der Parallelen« erst im Jahr 1804 an \textsc{Gauss} geschickt, dessen
Kritik der Brief vom 25. November 1804 enthält.

Die folgenden Notizen ergänzen die vorstehenden Äußerungen zu \textsc{Gauss}’
Untersuchungen über die ersten Gründe der
Geometrie.

\begin{quote}
\emph{»Plani possibilitatem demonstravi. [1797] Jul. 28. Gotting.«}
\end{quote}
und
\begin{quote}
\emph{»In principiis Geometriae egregios progressus fecimus.
Br[unovici 1799] Sept.«}
\end{quote}
Sie finden sich in dem \textls{Tagebuch}, das bereits auf S. 20 dieses Bandes
erwähnt wurde.

\begin{flushright}
\textsc{Stäckel.}
\end{flushright}

\medskip
\begin{center}
\rule{0.3\textwidth}{0.4pt}
\end{center}
